\scanpage{14}
\mspage{13}
\uline{Subdivision} d'une figure. \deleted{\unsure{Et|S}} On se reporte au cas-type. Soit une figure $F$, \unsure{formée|fermée} de $X$ et de \deleted{\unsure{l'ens|l'ens.}} de ses strates $X_i$.

Une \deleted{subdivision} figure $F'$ qui est une subdivision de $\deleted{F}\unsure{X|F}$ est \deleted{une} définie par les propriétés

\[
\text{a)}\quad |F'|=|F|
\]

b) la relation d'équiv. dans $X=|F|$ définie par $F'$ est plus fine que celle définie par \deleted{\unsure{$F$}} $F$.

La condition b) implique

b$'$) $\forall X_i$, strate (fermée) de $X$ pour $F$, $X_i$ est réunion de strates de $X$ pour $F'$, \deleted{\unsure{e|et}} et \inserted{elle} lui est $=$ équivalente.

Ainsi, les strates pour $F'$ qui sont $\subset X_i$, forment une \inserted{figure, qui est une} subdivision de $F_{X_i}$ (la figure \inserted{élémentaire} induite par $F$ sur $X_i$)

Mais on fera attention que cette figure, de support $X_i$, n'est pas néc. élémentaire (une subdivision d'une figure élém. peut n'être pas élémentaire)
